\section{Declarative Workflow Specification}
\label{sec:workflow}

% 1.5-2 pages.

\subsection{Capability versus composition}
% \AgentStepClass (cluster-scoped, admin): images, commands, executor, timeout,
%   evidence format, autofix, catalog dependencies. This is CAPABILITY.
% \AgentWorkflow (namespaced, team): which classes, dependency edges,
%   invalidation globs, candidates, budgets, selection. This is COMPOSITION ONLY.
% Security consequence: a tenant authoring a workflow cannot introduce an image.
% RBAC falls out: write on workflows, read-only on classes.

\subsection{Declaring dependencies and invalidation}
% after: cost ordering and fail-fast.
% invalidatedBy: the ledger's invalidation rules, as data rather than code.
% Include the minimal and the full example listings side by side.

\subsection{Resolution and version pinning}
% Resolution order: explicit ref, tag match, repo match, sole workflow, default.
% SECURITY: never resolve from ticket free text. Body and comments are writable
% by anyone with issue access, including the agent. Resolve from trusted
% structured fields only, once, at admission, before any agent code runs.
% Pin specHash: editing a workflow mid-flight would silently void ledger entries.

\subsection{Non-composable integrity requirements}
% \AgentWorkflowDefaults injects required integrity steps via mutating webhook.
% A tenant must not be able to compose away the suppression meta-gate.
% Forward-reference Section 9.
